% Generated by D:/tools/scratch_qdd/mmrec/render.py from data_a.py - edit the data, not this file.
\subsection{A. Astra 활용 --- 무엇을 맡기고, 한 번의 답은 얼마나 맞고, 돈은 얼마를 쓰나}\label{rec:A}
\noindent\textbf{한 줄:} Astra만 할 수 있는 일로 역할을 좁히고(§82), 탐침으로 흐름 상수를 정하고(§86), 한 호출 정확도를 재어 v2 프롬프트를 채택했지만 7--10 cm 어긋남을 못 알아채는 병목(0/13)이 남았다. Astra에 맞춘 단독 조종은 7/7 성공. 유료 호출은 이제 사용자 명시 승인 뒤에만 한다.

읽는 법: 각 항목은 사용자 말(빨강) $\rightarrow$ 해 봄 $\rightarrow$ 결과 $\rightarrow$ 그래서(결정) $\rightarrow$ 다음 순서다. 커밋 해시와 시각(KST)은 저장소 기록 그대로이고, 숫자는 결과 문서 값 그대로다.

\paragraph{A1. Astra를 쓰는 이유가 있어야 한다 $\rightarrow$ Astra 전용 일 J1--J6}
{\small 09-25 23:15 -- 23:38 KST}
\begin{itemize}
  \item \textbf{사용자} (ul 74, 09-25 23:15 KST): \intent{그 아스트라라는 존재를 상위 오케스트라로 쓰는 이유가 존재해야만해 얘가 자기 혼자서도 로봇조작을 할 수 있는놈이기때문에 아스트라를 너무 별일 아닌데도 활용을 해버리면, 그냥 더 작은 vlm쓰는게 맞음 최대한 엄청 최대한 능력치를 다써야해 아스트라를 쓸 수 밖에 없는 이유를 만들 수 있는 장점발휘를 할 수 있게 세팅을 해야함 좀 찾아서 어케 활용할지를 더 구체화해봐봐}
  \item \textbf{해 봄:} 조사 \texttt{docs/\allowbreak{}research/\allowbreak{}astra\_\allowbreak{}role\_\allowbreak{}2026-\allowbreak{}09-\allowbreak{}25.\allowbreak{}md}(925b0b0, 23:34): Astra만 할 수 있는 일, 작은 모델로 내릴 일, 호출 사다리, 필요성 실험 E-Astra-necessity 초안.
  \item \textbf{결과:} 강등 근거(정본 §82): V1h 세계 술어 0.949·0.44 ms 대 Astra 첫 토큰 약 3 s; FORTRESS(CoRL 2025) 즉석 GPT-4o 판정 0.64 대 사전 계산 경로 0.90.
  \item \textbf{사용자} (ul 76, 09-25 23:38 KST (질문: 하트비트·단계 경계 확인을 작은 모델로 내리고 Astra는 J1--J6에만 쓸지) 선택): \intent{바꾸기 (Recommended)}
  \item \textbf{사용자} (ul 76, 09-25 23:38 KST): \intent{논문에 로우 하이를 모두 가져가는 쪽이 맞겠네 성능비교두 그렇고}
  \item \textbf{사용자} (ul 76, 09-25 23:38 KST): \intent{지금 10만원 정도는 비용 한도 책정이 가능해}
  \item \textbf{그래서(결정):} 정본 §82(a1bb055, 23:38): Astra는 J1 과제 컴파일·J2 실패 진단과 복구안·J3 검사 술어 합성·J4 새 물체 이름·J5 저빈도 진행 감사·J6 경험 증류에만 쓴다. 5 s 하트비트는 J5로 대체, effort는 low·high를 모두 논문에 보고, 유료 한도 약 10만 원(4eb988d).
  \item \textbf{다음:} E-Astra-necessity(같은 자리에 로컬 VLM) 사전 등록 초안(2eea17b, 09-26 13:42)은 유료 실행 금지 상태로 남아 있다(A10).
\end{itemize}

\paragraph{A2. Astra가 움직임에도 개입 $\rightarrow$ 흐름 교정 설계, 호출은 한 번에 하나}
{\small 09-26 00:16 -- 02:18 KST (+ 14:14)}
\begin{itemize}
  \item \textbf{사용자} (ul 78, 09-26 00:16 KST): \intent{나는 그냥 아스트라 자체가 이 움직임에도 개입을 했으면 좋겠어 근데 딜레이가 문제니까 이걸 해결하면서 도 개입을 하는 방법을 찾는거야 최소화 하는 것도 방법이겠지 VLA가 한거를 감사하는건 아스트라의 능력중 극히 일부일 뿐인거구}
  \item \textbf{사용자} (ul 78, 00:18 KST (A: M4 느린 투표자 / B: 물체 기준 경로 / C: 조건부 프로그램)): \intent{A,B 더하는 방식이 조금 더 좋은 것 같긴하다}
  \item \textbf{사용자} (ul 83, 01:17 KST): \intent{로우로하자 그럼 이 로우와 하위 vla가 rtc로서 서로서로 보완하면 좋을듯 아스트라는 반응성이 느려서 그러지 vla는 방향성이 빠르지만 자기가 뭘하는지에대한 성능이 조금 떨어져서 그렇지}
  \item \textbf{사용자} (ul 83, 01:21 KST): \intent{이게 뭐랄까 엄청 촘촘하게 이루어질 수 있어야하는데 api를 직렬이 아니라 1초단위로 3개 불러오면 1초1초가 연속적인 방법이 될 수는 있잖어}
  \item \textbf{사용자} (ul 83, 01:33 KST): \intent{목표를 공간에 표현하는 것 조차 안하고 그냥 일반적 아스트라 조종법을 따르겠다는 얘기}
  \item \textbf{사용자} (ul 83, 01:33 KST): \intent{네, 일반 조종법만}
  \item \textbf{사용자} (ul 83, 01:40 KST): \intent{아스트라는 3개의 캠을 다봐야해 예를들어 헤드캠에서 봤을때는 원시때문에 물체를 잡을 때, 잡았다고 오탐하지 않게 우론목에는 안보이는 그런게 있는거지 그걸 해결할 수 이썽야하는거지}
  \item \textbf{해 봄:} 결합 설계 초안(2cdef56, 01:26): 손끝 목표 파이프라인, 1 Hz 계단식 Astra low, 두 층 M4, RTC 이음. 개정(ac165ab, 01:33): 공간 목표 표현 없이 일반 조종(continue/edit/stop + 평가). 설계 §11 부드러운 반영(00a0aad), §12 카메라 3대와 손목 근거(f71f407).
  \item \textbf{그래서(결정):} 정본 §82 보충 2(a8259df, 01:17): 실행 중 Astra effort = low, Astra와 VLA는 RTC식으로 서로 보완. high는 비교 실험·오프라인 일에서만.
  \item \textbf{사용자} (ul 86, 02:18 KST): \intent{그냥 아스트라 부르는거 그냥 1번씩만 하는걸로하자}
  \item \textbf{그래서(결정):} 정본 §84 보충 2(190af17, 02:18): 흐름 호출은 동시 요청 1개 --- 답(또는 시간 초과)이 오면 그때의 최신 관측으로 다음 요청. 1 s 계단식·동시 $\leq$ 3을 대체. `동시 1개'는 Claude 해석이며 사용자에게 알렸다(에피소드당 한 번이라는 뜻이면 정정).
  \item \textbf{사용자} (ul 96, 09-26 14:14 KST): \intent{주기로 아스트라 부르는건 끄는거?}
  \item \textbf{사용자} (ul 96, 09-26 14:14 KST): \intent{그 주기로 부르는건 켜두고.}
  \item \textbf{그래서(결정):} 흐름 호출은 켠 채 유지(17f7513). 새로움 게이트(E-NOV0)는 선택지 시험일 뿐이고 호출 방식 변경은 사용자 결정.
\end{itemize}

\paragraph{A3. E-Astra-motion 탐침 설계 --- 돈을 쓰기 전에 유익한가부터}
{\small 09-26 00:28 -- 03:27 KST}
\begin{itemize}
  \item \textbf{사용자} (ul 79, 00:28 KST): \intent{탐침 E-Astra-motion: 3D 경로 그리기 대 보고 조종하기 이거 한번 해보자용}
  \item \textbf{사용자} (ul 79, 00:31--00:37 KST): \intent{뎁스를가지고 한다고 하면 될 것 같아 뎁스가 있잖아 ai워커가}
  \item \textbf{사용자} (ul 79, 00:31--00:37 KST): \intent{아니다 뎁스를 쓰면 너무 범용성이 떨어질 것 같은 생각이드네 이미지 상만으로 어떻게 3D를 방향을 제대로 계획을 시킬 수 있을까...?}
  \item \textbf{사용자} (ul 79, 00:31--00:37 KST): \intent{그냥 뎁스라기보다는 그 가우시안 스플레팅인가? 아니면 포인트 클라우딩 포인트 클라우딩으이었던거 같은데 스테레오캠으로 하는 그거 그거를 어떻게 잘 해서 안되는건가? 공간 좌표 만들고 거기서 하는거.?}
  \item \textbf{해 봄:} 사전 등록 a62164d(02:39, 문서 시각 17:38:25Z, 유료 호출 전): 직렬 흐름, 요청 방식 F0/F1, 잡기 판정(카메라 구성), 부드러운 반영.
  \item \textbf{사용자} (ul 87, 02:41--02:45 KST): \intent{그 돈쓸떼는 정말 이게 유익한 실험이 되는지도 중요해}
  \item \textbf{사용자} (ul 87, 02:41--02:45 KST): \intent{고 나에게 승인은 안받아도되지만 자체적으로 저걸 해보라는거지}
  \item \textbf{사용자} (ul 87, 02:41--02:45 KST): \intent{그리고 만약 중간에 잘못되면 다시 새롭게 설계헤ㅐ서 낭비 없게 해야해}
  \item \textbf{사용자} (ul 87, 02:41--02:45 KST): \intent{매번 자체검사를 해야해 알겠지?}
  \item \textbf{결과:} 유료 일시 정지(5320dd2, 02:41). 옛 설계 G1은 13호출 204원(호출당 15.7원, 지연 p50 4.2 s). 무료 Qwen 사전 실행에서 카메라 3대일 때 `잡았다' 과잉 주장 12/14 $\rightarrow$ 원인은 `쥠'을 정의하지 않은 프롬프트, 정의 뒤 1/7(정본 §86).
  \item \textbf{그래서(결정):} 재범위(075c505, 02:43): 무료 Qwen 먼저, 유료는 G1(약 40장·4조건) + S2(F0·F1 각 3편) + G2(low 오답에만 high $\leq$ 8회), S1에서 Astra 뺌, 누적 하드 정지 15,000원; 관문마다 결함이면 멈추고 재설계. 자체 검사 규칙을 CLAUDE.md로(ed478e4, 02:45).
  \item \textbf{참고:} 등록 수정 1·2 커밋(35029e5, 03:27)은 유료 호출 재개보다 늦게 커밋되었다 --- 이탈로 기록(R7 25회차 E-D2).
\end{itemize}

\paragraph{A4. 탐침 결과 $\rightarrow$ 결합 상수 확정(§86)}
{\small 09-26 04:09 KST}
\begin{itemize}
  \item \textbf{결과:} 요청 방식: 뒤집힘율 0.594 $\rightarrow$ 0.381($-$36 \%, 등록 규칙 50 \% 감소 미달). 직렬 흐름(카메라 3대) 지연 p50 약 9--10 s·p95 12--16 s, 답 나이 p50 6.7--9.8 s --- 6 s 문턱이면 답의 70--100 \%가 버려진다.
  \item \textbf{결과:} 잡기 판정 40장(쥠 25/안 쥠 15): 머리만이면 uncertain 39/40, 손목 포함이면 오탐 0--1/15이지만 실제 쥠 13/25를 놓침(한 손가락 테두리 걸기 6/6 전부), high 표적 8장 중 2장만 고침; 덧그림 효과 없음(짝 3 대 3); `불확실 $\rightarrow$ continue' 규칙 발동 0회.
  \item \textbf{결과:} 부드러운 반영(개루프 재생): 저크 RMS 중앙 0.83 대 즉시 2.36 m/s$^3$, 최대 속도 0.036 대 0.095 m/s. Astra 단독($\leq$ 5 cm 수정, VLA 없음)은 물체 도달 0/6, Qwen 동기 0/20. 비용: 흐름 49.7--58.0원/호출, 파지 질문 16.8원, high 표적 27.9원, 로봇 1분당 약 350--500원, 캐시 적중 0, 탐침 총 6,933.4원(하드 정지 15,000원).
  \item \textbf{그래서(결정):} 정본 §86(3500f54): 요청 방식 = F0; 늦은 답 폐기 6 $\rightarrow$ 15 s(보충: 시간 초과 20 s); 파지 확정은 Astra 영상이 아니라 고유감각 T1·V1h(Astra 주장은 참고), 손목 영상은 Astra 입력에 유지; 확신 필드로 거르지 않고 두 답 합의·편향 속도 상한으로; 부드러운 반영 유지; 모든 판정 질문 프롬프트에 판정 대상의 조작적 정의. \texttt{CoupleParams}: F0·15 s·20 s·\texttt{latency\_\allowbreak{}init\_\allowbreak{}s} 9.3 s·\texttt{est\_\allowbreak{}text\_\allowbreak{}tokens} 1601(§86 보충 2).
  \item \textbf{다음:} Astra 단독 0/6은 뒤에 사용자가 `실험을 Astra에 맞게 안 한 탓'이라고 지적 $\rightarrow$ A8.
\end{itemize}

\paragraph{A5. 프롬프트 검진}
{\small 09-26 14:14 -- 15:51 KST}
\begin{itemize}
  \item \textbf{사용자} (ul 96, 14:14 KST): \intent{프롬프트 엔지니어링 등으로 문제가 생겼는지 검진하는게 필요해보이는데}
  \item \textbf{해 봄:} 모든 프롬프트(Astra 흐름·J1--J6·탐침 판정 질문·로컬 VLM 어댑터·VLA 직렬화 문장)를 정적 감사하고, 무료 Qwen3-VL 8B·4B로 문구·보기 순서·카메라 순서·정의를 흔들고, 유료 Astra 표본(사전 등록 5b62187, 14:39, 상한 3,000원·80 \% 정지).
  \item \textbf{결과:} 1b68a6a(15:16): \texttt{astra-\allowbreak{}couple@v1}에 판정 정의·축 문구가 없고, 그리지 않은 화살표를 범례가 설명(Qwen 8B 인용 G 27/33), 확정 화살표 = 결정 토큰 중심값. Qwen 문구 뒤집힘 8B 6--17/40·4B 최대 28/40(같은 문구 반복 0/40), 방향 질문은 우연 수준. Astra 잡기 문구 뒤집힘 3/30(강건), 축 문구 18/30 대 17/30(가를 수 없음, x 부호 오류 = 원근). VLA \texttt{mag\_\allowbreak{}coarse}·\texttt{dir\_\allowbreak{}xy} 라벨 뜻 $\neq$ 질문 문구(청크 일치 0.375·0.626) $\rightarrow$ E-SR0·E-SR1b의 부분 원인. 탐침 Astra 답의 도착 시점 방향 F0 25/40·F1 20/24. 2,093.9원, 약 0.8 GPU-h.
  \item \textbf{그래서(결정):} 운영 프롬프트는 그대로 두고(판본 기록 누락 1건만 고침) 발견은 새 판에: Astra 쪽은 \texttt{astra-\allowbreak{}couple@v2}(A6), VLA 쪽은 PH-A1(b931f75, 15:51) --- 라벨은 그대로, 융합 VLA 문구를 라벨 정의대로. 로컬 VLM 이미지 상한 2 $\rightarrow$ 3(§59 보충, a913020).
\end{itemize}

\paragraph{A6. Astra 한 번의 명령은 정확해야 한다 $\rightarrow$ E-ACC 1단계, v2 채택}
{\small 09-26 14:24 -- 19:41 KST}
\begin{itemize}
  \item \textbf{사용자} (ul 97, 14:24 KST): \intent{아스트라가 다음 답변을 할때까지 아무것도 못하는 거니까. 아스트라가 한번 명령을 내릴때 좀 정확히 해야함}
  \item \textbf{사용자} (ul 104, 15:24 KST): \intent{아스트라를 최대 성능이 나게 만들어놓는것이 광장히 중요함}
  \item \textbf{사용자} (ul 104, 15:24 KST): \intent{vla는 아스트라가 옳다는 전제로 움직이기때문}
  \item \textbf{그래서(결정):} 정본 §94(64b0a45, 15:24): Astra 호출 단위 정확도 최대화가 E-Couple의 선행 조건, E-ACC를 먼저.
  \item \textbf{해 봄:} 사전 등록 e4f8891(15:52): Isaac 벤치(R2 오라클 계획기 $\times$ 0.5 속도를 대역 정책으로), 스냅숏 종류 on·off\_a(7--10 cm 엉뚱한 곳)·off\_b(엉뚱한 물체)·off\_c(엉뚱한 놓을 곳), 팔 v1·v2·v2cp·v2\_ax, 채점 = 도착 시각(스냅숏 + 9.3 s) 기준. 벤치 42편. 무료 Qwen 선별 G-fmt 실패 $\rightarrow$ 변경 2(eb42189)로 형식만 고침.
  \item \textbf{결과:} 사고: P1의 31번째 호출부터 OpenAI 크레딧 소진(\texttt{insufficient\_\allowbreak{}quota}) --- 빈 답 37개, 과금 1,352.7원(on 30호출만). 클라이언트가 스트림 \texttt{error}·\texttt{response.\allowbreak{}failed}를 오류로 올리지 않았다(함정 P108, cfaf140 16:45).
  \item \textbf{그래서(결정):} 변경 3·4(a09c5ed, 19:05): 새 장부(과금 이월), quota 즉시 정지·빈 답 3연속 정지, off 20장 $\times$ v1·v2·v2\_ax로 재개, v2cp는 뺌.
  \item \textbf{결과:} 2a7b9ae(19:15): 무효 0/90. 수정(edit)을 냈을 때 방향은 13/13 맞음(도착 각 중앙 5--9$^\circ$). 그러나 off\_a·off\_c는 세 팔 모두 0/13(`progressing / aligned'). off\_b는 v1 1/7·v2 5/7·v2\_ax 7/7. R1 v2 $-$ v1: M1 +0.20 [+0.05, +0.40] $\rightarrow$ v2 채택. R2$'$ 축 덧그림: +0.10 [0.00, +0.25] $\rightarrow$ 미정, 기본 끔. 구간 now 정답 v2 26/30. 원/호출 v1 36.4·v2 49.1·v2\_ax 51.3.
  \item \textbf{그래서(결정):} \texttt{astra-\allowbreak{}couple@v2}(8c9064a, 19:41, 결합 Task 18) = E-ACC 채택 문장을 바이트 그대로(판본 83fa03a5de19): 축 줄, 정의(쥠 = 두 패드가 닫혀 물체에 붙음), 실제로 그린 요소만 범례, \texttt{segment\{now, do, next\}}, \texttt{valid\_\allowbreak{}until}, \texttt{since\_\allowbreak{}last\_\allowbreak{}request}. AxisGuide·카메라 자세 줄은 기본 끔.
\end{itemize}

\paragraph{A7. 알아채기 병목 --- E-ACC 2단계와 B$'$}
{\small 09-26 19:18 -- 19:34 KST}
\begin{itemize}
  \item \textbf{해 봄:} 변경 5(0dd95ad, 19:18): off\_a·off\_c 13장에서 (A) effort medium, (B) 목표 확인 문장(on 5장으로 괜한 수정도 확인).
  \item \textbf{결과:} 39ae001(19:28): medium 1/13, 지연 p95 17.6 s $\rightarrow$ 권고 안 함(low 유지). 목표 확인 7/13(+0.54 [+0.23, +0.77]; off\_a 6/7·off\_c 1/6)이지만 계획대로 상태 5장 중 괜한 수정 2장 $\rightarrow$ S2B 불채택. 괜한 수정의 원인은 `도착 예상 손끝' 외삽(0.5 s 속도 $\times$ 9.3 s)이 6--9 cm 튄 것(P113).
  \item \textbf{해 봄:} 변경 6(47e6ff6, 19:29): B$'$ --- 목표 확인의 기준점을 `지금 손끝 + 확정 화살표'로.
  \item \textbf{결과:} f54cc88(19:34): 괜한 수정은 사라졌지만(on 5/5) 알아채기 3/13(+0.23 [0.0, +0.46]), off\_c 0/6 $\rightarrow$ S2B$'$ 불채택. E-ACC 합계 6,892.2원(상한 8,000원).
  \item \textbf{그래서(결정):} 두 수단 모두 v2에 넣지 않는다. 대신 결합의 예상 상태 기본을 외삽이 아니라 `지금 손끝 + 실행 청크 변위'로(§91 구현, Task 19, 20d735c 20:38). 남은 병목: Astra는 7--10 cm 어긋난 목표·놓을 곳을 알아채지 못한다(0/13) --- VLA는 Astra가 옳다고 보고 움직이므로(§94) 결합의 가장 큰 오류 전파 경로로 기록(§90 구현 기록 3).
  \item \textbf{다음:} 보고만 한 후보: 외삽을 고친 뒤 B 다시 재기, 로봇 좌표 평면도(APC). 둘 다 새 등록과 상한 증액 검토가 필요해 실행하지 않았다.
\end{itemize}

\paragraph{A8. Astra 단독도 꽤 할 것 $\rightarrow$ Astra에 맞춘 조건에서 4/4}
{\small 09-26 15:18 -- 16:38 KST}
\begin{itemize}
  \item \textbf{사용자} (ul 103, 15:18 KST): \intent{아스트라 혼자조종도 꽤 할거야. 나는 vla와 아스트라의 하이브리드가 그리 옳다고는 생각안해}
  \item \textbf{사용자} (ul 103, 15:19 KST): \intent{3이 맞음 근데 아스트라가 0이 나온건 너가 최적으로 아스트라에 맞게 실험을 안했다는 얘기기도해}
  \item \textbf{그래서(결정):} 정본 §93(ed06d8a, 15:19): 하이브리드 유지. 탐침 0/6은 Astra에 불리한 조건(v1 프롬프트 결함, 수정 $\leq$ 5 cm, 90 s·30호출, 기다리지 않는 흐름)의 값이라 Astra 실력으로 인용하지 않는다. Astra 단독 팔은 Astra에 최적으로 다시 설계. 함정: 불리한 조건의 기준선은 비교가 아니다(P101).
  \item \textbf{해 봄:} 사전 등록 50336ce(16:07, 변경 1 포함, 첫 유료 호출 전): 동기(로봇이 답을 기다림), \texttt{astra-\allowbreak{}solo@v2}(정의·축·카메라 자세·그린 요소만 범례·측정 이력), 명령 eef 절대 목표 + edit $\leq$ 10 cm + gripper + stop, 최소 저크 실행기(VLA 없음), 머리 영상에 탁자면 로봇 좌표 격자(5 cm) + TCP 낙하선, 편당 40호출·180 s. 무료 단계: 참값 모델 6/6, Qwen3-VL-8B 0/5(인식 한계 --- 배선 결함 2개는 유료 전에 고침).
  \item \textbf{결과:} cd80896(16:38): Astra low 4/4(mug\_tray DEV standard 0·1, dr 0·1). 편당 9--17호출(평균 11.0), 성공까지 움직임 중앙 16.35 s, 벽시계 74--215 s, 성공 1건당 591.5원, 무효 0/44. dr 0은 첫 닫기가 49 mm 빗나간 것을 패드 간격으로 알아채고 다시 잡아 성공. 접근 목표 xy 오차(편 중앙의 중앙) 14.1 mm. effort 짝 12개 중 동률 아닌 짝 5 < 8 $\rightarrow$ undecided(medium 4승). 3,080.9원.
  \item \textbf{그래서(결정):} D1 ADOPT $\rightarrow$ 전체 팔은 \texttt{astra-\allowbreak{}solo@v2} 그대로; D2 $\rightarrow$ low 기본 + medium은 판정 밖 부분집합. 해석: E-Couple 사전 실행의 VLA 단독 0/12와 함께 보면 `Astra 단독(동기)'이 강한 기준선일 수 있다(§92 P3와 함께 봐야 함).
\end{itemize}

\paragraph{A9. Astra 단독 전체 팔 --- 크레딧 소진, 비용 지적, 멈춤}
{\small 09-26 18:42 -- 19:21 KST}
\begin{itemize}
  \item \textbf{해 봄:} 등록 변경 2(d0d2ae7, 18:42): AS 팔을 §92 P3 기준선으로 --- 40편(DEV 0--19 $\times$ standard/dr), 하드 정지 25,000원; 러너 가드(\texttt{insufficient\_\allowbreak{}quota} 즉시 치명, 빈/오류 답 3연속 정지). OpenAI 크레딧 소진으로 처음에는 보류(HOLD).
  \item \textbf{결과:} 19:00 KST(10:00 UTC)에 시작해 3편 모두 성공(9·9·9호출, 460·456·443원), 장부 1,359.0원에서 메인이 멈춤.
  \item \textbf{사용자 말(user-log 번호 없음)} (메인 전달, 등록 변경 3에 인용): ``한 번에 23,000원 쓰는 건좀 아닌 것 같은데 좀 아껴서 하지 제대로 된 결과 나오게 아껴서 하지''
  \item \textbf{그래서(결정):} 등록 변경 3(11cbc39, 19:13): 단계식 재설계 --- 1단계 = DEV 0--5 $\times$ standard/dr 12편(이미 3편), 1단계 하드 정지 8,000원; 러너 조기 종료(20호출 또는 움직임 60 s, 프롬프트 바이트는 그대로)로 실패 편 최대 약 2,100 $\rightarrow$ 1,050원; 순차 규칙($\geq$ 10/12 강함·$\leq$ 2/12 약함에서 멈춤). 비용 분해: 입력 약 63 \%, 캐시 적중 0/80.
  \item \textbf{진행 중·대기:} \tentative{1단계 나머지 9편은 결과 문서가 없다(실행 기록 없음). 이후 규칙(ul 114)으로 재개에는 사용자 명시 승인이 필요하다. `Astra 단독 7/7' = 파일럿 4/4 + 전체 팔 앞 3/3.}
  \item \textbf{참고:} Astra 단독 실제 편 7개 영상을 파드 /data에 저장(4d54506, 19:21) --- 영상 규칙의 첫 적용(F8).
\end{itemize}

\paragraph{A10. 비용 사건과 유료 승인 규칙}
{\small 09-26 16:45 -- 09-27 00:51 KST}
\begin{itemize}
  \item \textbf{결과:} 사건 목록: (1) E-ACC P1 크레딧 소진 --- 빈 답 37개, 장부 예약 5,043원은 미과금, 과금 1,352.7원(P108); (2) AS 전체 팔 크레딧 HOLD; (3) AS 전체 팔 3편 뒤 비용 지적으로 멈춤(A9); (4) E-OpenVLM 대리의 gpt-5-mini가 7호출(23.3원) 뒤 사용자 지시로 멈춤(결과 \texttt{open\_\allowbreak{}vlm\_\allowbreak{}solo.\allowbreak{}md} 1절, B3).
  \item \textbf{결과:} 지출(책 05): 확정 누적 19,000.4원 = 탐침 6,933.4 + 프롬프트 검진 2,093.9 + Astra 단독 파일럿 3,080.9 + E-ACC 6,892.2; 여기에 E-OpenVLM 대리 1,322.4원을 따로 더한다(책 05 표기). AS 전체 팔 앞 3편 1,359.0원은 그 장부(등록 변경 3).
  \item \textbf{사용자} (ul 114, 09-26 21:55 KST): \intent{돈드는건 내 허락맡고 하기}
  \item \textbf{그래서(결정):} user-log 114(2882208, 21:56): 유료 호출(Astra·OpenAI 등)은 목적·호출 수·예상 비용을 먼저 제안하고 사용자 명시 허락 뒤에만 --- ul 87의 `승인 불필요'를 대체, 자체 검사 의무는 그대로. 가드·등록 문구(99b3b58, 22:23, Task 26), 모든 실제 Astra 클라이언트는 승인 근거 없이 거부·\texttt{run\_\allowbreak{}r5} 기본 mock(fddefd4, 23:59, §85 후속 2), CLAUDE.md 줄(02de6cb, 09-27 00:51). 앞서 \texttt{insufficient\_\allowbreak{}quota} 치명 정지·미응답 요청 별도 보고(688992c, 19:14, Task 21).
  \item \textbf{다음:} 유료 E-Couple(상한 25,000원)·E-Astra-necessity·AS 1단계·자동 맥락 A2 유료 절제·E-AT는 모두 사용자 승인 전 실행 금지.
\end{itemize}
